% !TEX program = xelatex
\documentclass[a4paper,12pt]{report}

\usepackage[margin=2.5cm]{geometry}
\usepackage{fontspec}
\usepackage{xeCJK}
\usepackage{graphicx}
\usepackage{float}
\usepackage{booktabs}
\usepackage{multirow}
\usepackage{array}
\usepackage{longtable}
\usepackage{listings}
\usepackage{xcolor}
\usepackage{hyperref}
\usepackage{enumitem}
\usepackage{setspace}
\usepackage{fancyhdr}
\usepackage{titlesec}

\setCJKmainfont[AutoFakeBold=2]{Microsoft JhengHei}
\setCJKsansfont[AutoFakeBold=2]{Microsoft JhengHei}
\setCJKmonofont[AutoFakeBold=2]{Microsoft JhengHei}
\setmainfont{Times New Roman}

\XeTeXlinebreaklocale "zh"
\XeTeXlinebreakskip = 0pt plus 1pt

\definecolor{codebg}{RGB}{245,247,249}
\definecolor{lineblue}{RGB}{0,122,204}

\lstset{
    language=Python,
    backgroundcolor=\color{codebg},
    basicstyle=\ttfamily\small,
    keywordstyle=\color{blue},
    commentstyle=\color{gray},
    stringstyle=\color{purple},
    showstringspaces=false,
    breaklines=true,
    frame=single,
    numbers=left,
    numberstyle=\tiny,
    stepnumber=1,
    rulecolor=\color{black},
    captionpos=b,
    keepspaces=true,
    columns=fullflexible,
    xleftmargin=1em,
    tabsize=4
}

\renewcommand{\contentsname}{目錄}
\renewcommand{\listfigurename}{圖目錄}
\renewcommand{\listtablename}{表目錄}
\renewcommand{\figurename}{圖}
\renewcommand{\tablename}{表}
\renewcommand{\lstlistingname}{程式}
\renewcommand{\lstlistlistingname}{程式目錄}
\renewcommand{\chaptername}{第}
\renewcommand{\thechapter}{\arabic{chapter}章}
\renewcommand{\thesection}{\arabic{chapter}.\arabic{section}}
\renewcommand{\thesubsection}{\thesection.\arabic{subsection}}

\title{LINE Messaging API 註冊申請與應用實作}
\author{作者姓名}
\date{\today}

\begin{document}
\maketitle

\begin{abstract}
本文件整理 LINE Messaging API 的註冊、設定與實務應用流程，內容涵蓋 LINE Developers 帳號建立、Channel 設定、Channel Secret 及 Access Token 的取得方式、Webhook 設定與測試方法，以及使用 Python 範例程式接收與回覆訊息的實作流程。透過本文件，讀者可快速建立對 LINE Messaging API 的基本理解，並具備在專案中接入 LINE 訊息服務的基礎能力。
\end{abstract}

\tableofcontents
\listoffigures
\listoftables
\lstlistoflistings

\chapter{LINE Messaging API 簡介}
\section{LINE Messaging API 的定位與應用場景}

LINE Messaging API 是 LINE 官方提供的一套訊息通訊能力，讓開發者可透過程式自動收發訊息、推播通知、進行互動式問答，以及實作客服與活動通知等功能。它不只是單純的聊天工具，而是一個可嵌入商業流程的開發介面。企業與團隊常利用它在客服、會員通知、購物提醒、活動回饋、內部通知與自動化流程中實際應用。

\subsection{API 的主要價值}

在一般商業場景中，傳統的人工客服與簡訊通知往往存在成本高、回應速度慢、訊息彙整不易等問題。透過 LINE Messaging API，開發者可以將 LINE 變成企業與使用者之間的數位通道，讓系統自動回覆訊息、傳遞訂單資訊、發送預約提醒或通知會員活動。

例如，電商平台可以在訂單狀態更新時，自動推送「已付款」「已出貨」等訊息；醫療或教育系統可以發送預約提醒與報名結果；客服系統可在使用者輸入關鍵字後由機器人回覆對應回答。這些場景都能透過 LINE Messaging API 來實作，而不需要建立額外的通訊平台。

\subsection{常見的 LINE Bot 應用}

LINE Bot 的使用方式相當多元，常見的類型包括：
\begin{itemize}[leftmargin=2em]
    \item 會員通知：發送免責說明、帳戶提醒、點數異動通知。
    \item 客服自動回覆：透過關鍵字或選單進行 FAQ 問答。
    \item 活動互動：使用者可透過 LINE 回覆訊息參與抽獎或報名。
    \item 內部協作：將 LINE 作為通知中心，降低各部門溝通成本。
    \item 數據收集：接收使用者回覆內容，進行後續分析與行銷用途。
\end{itemize}

\subsection{LINE Messaging API 與一般 API 的差異}

LINE Messaging API 的設計重點在「訊息互動」與「事件驅動」。它與一般資料 API 的差異在於：訊息是由事件觸發，而不是單純讀取一段資料。當使用者發送訊息、點擊按鈕或加入好友後，系統會產生事件，開發者就能接收這些事件並做後續處理。

這種設計使 LINE Messaging API 特別適合做即時互動。舉例來說，當使用者在聊天室中輸入文字時，Webhook 會自動通知伺服器，伺服器再依照規則做回應。這樣的流程使互動能保持近乎即時。

\begin{figure}[H]
    \centering
    \fbox{\rule{0pt}{3.2cm}\rule{9cm}{0pt}}
    \caption{LINE Messaging API 通訊流程示意圖}
\end{figure}

\subsection{技術觀點的理解}

從技術架構來看，LINE Messaging API 由三個核心元素組成：Channel、Token 以及 Webhook。Channel 是應用程式的身份資料，Token 是驗證與授權憑證，Webhook 則讓伺服器在使用者事件發生時取得即時訊息。這三者配合使用後，開發者才能完成訊息收發與安全驗證。

也就是說，若想把 LINE 接進專案，最重要的並不是會撰寫單一程式，而是理解這一整套訊息流程：先建立 Channel、再取得 Secret 與 Token、再透過 Webhook 接收事件，最後依據事件內容回覆訊息或觸發自動化流程。這樣的設計使 LINE API 不僅能傳送訊息，也能成為自動化業務運作平台。

\section{LINE Messaging API 的設計原則}

LINE 的訊息平台強調安全、穩定與實作簡潔。它要求開發者妥善管理 Channel Secret、Access Token 與訊息驗證邏輯，否則可能導致伺服器被偽造請求侵入，或產生未預期的訊息發送行為。良好的設計通常會遵循以下幾個原則：

\begin{itemize}[leftmargin=2em]
    \item 不要直接將 Secret 或 Token 存放在前端程式碼中。
    \item 使用環境變數或外部設定檔管理敏感資訊。
    \item Webhook 需驗證簽名與來源，避免被任意請求觸發。
    \item 回應訊息前先進行錯誤處理與資料驗證。
    \item 以服務端為核心，將 LINE 互動邏輯與業務邏輯分離。
\end{itemize}

\begin{table}[H]
    \centering
    \caption{LINE Messaging API 典型使用情境}
    \begin{tabular}{|p{4cm}|p{8cm}|}
        \hline
        \textbf{使用情境} & \textbf{說明} \\
        \hline
        客服問答系統 & 使用者輸入常見問題時，系統自動判斷並回覆答案。 \\
        \hline
        訂單與物流通知 & 例如訂單成立、出貨、配送、取貨成功等自動通知。 \\
        \hline
        活動報名 & 使用者可透過訊息參與活動並立即收到確認回覆。 \\
        \hline
        會員提醒 & 系統發送會員生日、購物回饋、促銷提醒等內容。 \\
        \hline
    \end{tabular}
\end{table}

從實務角度來看，LINE Messaging API 不是單一功能，而是一個可擴充的通訊基礎設施。只要掌握消息事件、回應機制與安全設定，開發者便可將它擴展到客服、個人化推播、報名管理與自動化通知等多種場景。


\chapter{LINE Developers 帳號與 Channel 建立}
\section{LINE Developers 帳號註冊與 Channel 建立}

想要使用 LINE Messaging API，開發者首先需要準備一個 LINE Developers 帳號。若尚未擁有帳號，需先登入 LINE 官方帳號並進入 LINE Developers 官網。通常，以個人或團隊的官方身分建立帳戶後，即可建立自己的專案。

\subsection{註冊步驟}

登入 LINE Developers 後，進入「Create a provider」或「建立 Provider」頁面。Provider 是一個開發者專案群組，用於管理多個 LINE Bot 或其他 LINE 服務。建立 Provider 後，接著點選「Create a channel」建立新的 Channel。Channel 類型通常可選擇 Messaging API，這是用於發送與接收訊息的主要設定項目。

\begin{enumerate}[leftmargin=2em]
    \item 登入 LINE Developers 平台。
    \item 建立 Provider，命名公司或專案名稱。
    \item 進入 Provider 內部選擇「Create a channel」。
    \item 選擇「Messaging API」Channel 類型。
    \item 填寫 Channel 名稱、描述與相關資訊。
    \item 完成設定後，建立成功並取得 Channel ID。
\end{enumerate}

\subsection{Channel 設定說明}

當 Channel 建立成功後，開發者進入 Channel 設定頁面，這裡會看到關鍵資訊，例如：
\begin{itemize}[leftmargin=2em]
    \item Channel Secret
    \item Access Token
    \item Webhook URL
    \item 允許的回覆與推播設定
    \item 監控與錯誤資訊頁面
\end{itemize}

這些設定決定了訊息能否被伺服器正確接收與回覆。舉例來說，若 Webhook URL 未正確設定，當使用者傳送訊息時，系統不會有任何事件回應。若 Access Token 未正確配置，則伺服器也無法成功發送訊息。

\subsection{Line Bot 的建立概念}

LINE Bot 是一個透過 LINE Messaging API 運作的智慧型聊天功能，通常由後端服務來負責處理邏輯。商業團隊可以在自己的系統內部管理資料，並將 LINE 作為人機介面。藉由這種方式，開發者可以用最小的前端成本，建立一個具備對話能力的服務入口。

在正式上線前，開發者通常會先使用測試帳號或用戶端進行功能驗證，再將設定同步到正式環境。這樣可以降低錯誤率，並讓機器人更容易在正式動線中穩定運作。

\subsection{Channel 相關注意事項}

建立 Channel 時需要注意幾個實務層面：
\begin{itemize}[leftmargin=2em]
    \item 不要把 Channel Secret 或 Access Token 直接寫死在前端 HTML 中。
    \item 若使用 Git 進行版本控，請將設定檔加入 .gitignore。
    \item 若測試環境與正式環境分離，最好使用不同的 Channel 或各自的環境變數。
    \item 在正式上線前，務必確認 Webhook 網域可被公開存取。
\end{itemize}

\begin{table}[H]
    \centering
    \caption{LINE Developers 建立 Channel 的主要步驟}
    \begin{tabular}{|p{4cm}|p{8cm}|}
        \hline
        \textbf{步驟} & \textbf{說明} \\
        \hline
        建立 Provider & 管理多個 Channel 與專案資源。 \\
        \hline
        建立 Channel & 選擇 Messaging API 類型並命名專案。 \\
        \hline
        設定 Webhook & 提供公開 URL 讓 LINE 可回傳事件資訊。 \\
        \hline
        產生 Token & 取得 Access Token 以便後端發送訊息。 \\
        \hline
    \end{tabular}
\end{table}

在專案完成後，LINE Bot 可以用於推播通知、客服自動回覆、預約提醒甚至會員互動等多種功能。這些特性使 LINE Messaging API 成為企業與開發者最常用的互動通道之一。


\chapter{Channel Secret、Access Token 與 Webhook 設定}
\section{Channel Secret、Access Token 與 Webhook 設定}

LINE Messaging API 的安全與權限是整個串接流程中的核心。若沒有正確設定 Channel Secret 與 Access Token，即使已建立 Channel，也無法安全地接收與回應使用者訊息。這章節整理這些關鍵設定與實務注意事項。

\subsection{Channel Secret}

Channel Secret 是 LINE 平台提供給開發者的機密資訊，用於識別 Channel 身分。它通常在 LINE Developers 後台顯示，可在 Channel 基本設定中查看。這個值應該被視為高度敏感的憑證，應妥善保存，不要以明碼方式直接寫入應用程式碼。

Channel Secret 常見於兩種情境：
\begin{enumerate}[leftmargin=2em]
    \item Webhook 驗證與請求來源確認。
    \item 需要驗證該請求是否來自正確的 LINE Channel。
\end{enumerate}

\subsection{Access Token}

Access Token 是讓伺服器能發送訊息或呼叫 API 的權限憑證。通常分為「一般權限 Token」與「長期 Token」。對新手而言，先理解它的作用最重要：它代表「伺服器具備對該 Channel 的操作權限」。

使用 Access Token 的常見操作包括：
\begin{itemize}[leftmargin=2em]
    \item 傳送文字訊息
    \item 回覆使用者訊息
    \item 發送多媒體訊息
    \item 推播通知給特定使用者或群組
\end{itemize}

如果 Access Token 過期或被重置，原本的程式可能無法再發送訊息，因此系統設計時最好透過環境變數分開管理，並處理例外情況。

\subsection{Webhook 設定}

Webhook 是 LINE 事件通知入口。當使用者傳送訊息、點擊按鈕或發生特定事件時，LINE 伺服器會 POST 請求到預先設定的 Webhook URL。開發者可以在伺服器中接收這些請求，分析事件內容，進而回覆訊息。

Webhook 的常見架構如下：
\begin{itemize}[leftmargin=2em]
    \item 伺服器公開 URL，例如 https://example.com/callback
    \item 使用 HTTPS 讓 LINE 可安全訪問
    \item 驗證簽名，確保事件資料真的來自 LINE
    \item 解析 JSON，提取 replyToken、訊息內容或使用者資訊
\end{itemize}

若使用本機開發環境，例如測試伺服器，通常需要透過 ngrok 把本地端口轉發到外網，以讓 LINE 可以訪問。ngrok 的使用方式簡單，開發者只需執行 ngrok http 8000，便可取得一個公開 URL。

\subsection{使用 ngrok 的常見作法}

在開發階段，開發者不一定能直接將應用程式部署到正式伺服器，因此會使用 ngrok 來產生測試入口。這有助於快速驗證 Webhook 是否可用，並確保事件回發正常。這樣的做法在機器人開發中相當常見，因為它能把本地端點對外暴露，讓 LINE 服務可直接訪問。

\subsection{安全實務}

LINE API 的安全設定不只是把設定填好而已，還需要遵守一套良好的開發實務：
\begin{enumerate}[leftmargin=2em]
    \item 公開 Webhook 前，先使用測試帳號驗證。
    \item 使用 HTTPS，不要使用無保護的 HTTP。
    \item 不要洩漏 Channel Secret 或 Access Token。
    \item 驗證訊息來源與數據完整性。
    \item 啟用日誌紀錄，方便排查錯誤。
\end{enumerate}

\begin{figure}[H]
    \centering
    \fbox{\rule{0pt}{2.5cm}\rule{12cm}{0pt}}
    \caption{LINE Messaging API 設定流程總覽}
\end{figure}

整體設定流程可概括為：建立 Provider → 建立 Channel → 取得 Secret 與 Token → 設定 Webhook → 驗證事件 → 接收與回覆訊息。這個流程是所有 LINE Bot 開發的基礎，也是後續業務整合的起點。


\chapter{LINE 訊息接收與回覆範例}
\section{LINE 訊息接收與回覆範例}

本節示範一個使用 Python Flask 建立 Webhook 伺服器的簡化範例，當使用者在 LINE 上傳送訊息時，伺服器將接收事件並回覆相同內容。這個例子可作為理解 LINE Messaging API 的進入點，也是一種最常見的初學範例。

\subsection{範例程式：接收 LINE Webhook 並回覆訊息}

\begin{lstlisting}[language=Python, caption={Flask Webhook 範例：接收訊息並回覆}, label={lst:webhook}]
from flask import Flask, request, jsonify
import os

app = Flask(__name__)

CHANNEL_SECRET = os.getenv("LINE_CHANNEL_SECRET", "your_channel_secret")
ACCESS_TOKEN = os.getenv("LINE_ACCESS_TOKEN", "your_access_token")

@app.route("/callback", methods=["POST"])
def callback():
    data = request.get_json()
    events = data.get("events", [])

    for event in events:
        if event.get("type") == "message" and event["message"]["type"] == "text":
            user_text = event["message"]["text"]
            reply_token = event["replyToken"]

            reply_message = {
                "replyToken": reply_token,
                "messages": [
                    {
                        "type": "text",
                        "text": f"你剛剛傳了：{user_text}\n我已收到你的訊息。"
                    }
                ]
            }

            print("收到訊息:", user_text)

    return jsonify({"status": "ok"})

if __name__ == "__main__":
    app.run(host="0.0.0.0", port=8000)
\end{lstlisting}

\subsection{程式說明}

上述範例有幾個關鍵點值得注意：
\begin{enumerate}[leftmargin=2em]
    \item 使用 Flask 建立 HTTP API，接收 LINE 的 POST 請求。
    \item 從 request.get\_json() 取得事件資料。
    \item 判斷 event 的 type 是否為 message。
    \item 使用 message["text"] 取得使用者輸入內容。
    \item 使用 replyToken 回覆使用者訊息。
\end{enumerate}

這段程式能夠完成最基本的雙向互動：使用者傳訊息，伺服器接收，並回傳相同或延伸資訊。實際專案中，開發者可以再加上關鍵字判斷、資料庫查詢、客服選單與群發通知功能。

\subsection{進一步擴充應用}

除了簡單地回覆文字之外，LINE Messaging API 還支援多種訊息型別，例如：
\begin{itemize}[leftmargin=2em]
    \item 文字訊息
    \item 圖片訊息
    \item 地圖與位置訊息
    \item 輪播訊息
    \item 按鈕模板
    \item 快速回覆
\end{itemize}

此外，LINE Bot 還可與資料庫整合，將使用者輸入內容儲存到後端系統，或實作應用程式內的客服自動回覆流程。這使其非常適合用於 CRM、會員服務與自動化問答平台。

\subsection{錯誤排查與測試建議}

在整合 LINE Messaging API 時，常見錯誤來源包括：
\begin{itemize}[leftmargin=2em]
    \item Webhook URL 尚未設定或未公開
    \item Access Token 無效或失效
    \item 伺服器沒有正確回應 HTTP 200
    \item JSON 格式不符合預期
    \item Channel Secret 與實際設定不一致
\end{itemize}

因此，建議開發者在測試階段搭配 nginx、ngrok 或測試伺服器，確保 LINE 事件透過公開網址送達，並在本機或雲端環境中逐步驗證每一個訊息事件。

\subsection{系統流程總結}

LINE Messaging API 的完整流程可以概括為：
\begin{enumerate}[leftmargin=2em]
    \item 建立 LINE Developers Provider 與 Messaging API Channel
    \item 取得 Channel Secret 與 Access Token
    \item 設定 Webhook URL
    \item 伺服器接收使用者事件
    \item 判斷事件類型與內容
    \item 回覆訊息或觸發其他流程
\end{enumerate}

這樣的流程讓 LINE 變成高度可擴充的互動入口，適合多種企業與個人專案使用。


\chapter{本文件 LaTeX 原始碼}
\section{程式碼清單}
本章節直接展示本文件的主要原始碼，供讀者檢視整份 LaTeX 文件的結構與設定。

\lstinputlisting[
    language=TeX,
    caption={main.tex 原始碼},
    label={lst:mainSource},
    breaklines=true,
    numbers=left
]{main.tex}

\end{document}